\chapter{Prosperity Algorithmic Rounds: Strategies Mapped to Theory}
\label{ch:algo_rounds}
\chaptermark{Algorithmic rounds: strategies mapped to theory}

Prosperity products fall into a small set of recurring archetypes.
Recognise the archetype and you know which chapter to open.  This
chapter maps archetype~$\to$~theory~$\to$~strategy skeleton~$\to$~pitfalls
so a new product can be diagnosed in under five minutes.

\begin{backgroundbox}[Prerequisites]
\begin{itemize}
  \item \Cref{ch:market_making}: Avellaneda--Stoikov, inventory risk,
  fair-value quoting.
  \item \Cref{ch:mean_reversion_ou}: Ornstein--Uhlenbeck, z-scores.
  \item \Cref{ch:pairs_cointegration}: cointegration, spreads,
  hedge ratios.
  \item \Cref{ch:basket_etf_arb}: basket/ETF stat-arb, premium
  tracking.
  \item \Cref{ch:momentum}: trend signals, EMA, regimes.
  \item \Cref{ch:arbitrage_zoo}: triangular, locational, calendar,
  cash-and-carry arb.
  \item \Cref{ch:black_scholes,ch:greeks_hedging,ch:vol_surface}:
  Black--Scholes, Greeks, delta hedging, IV, vol smile.
  \item \Cref{ch:which_strategy_when}: four-layer observables$\to$strategy
  pipeline.
  \item \Cref{ch:prosperity_setup}: Prosperity API, timestep model,
  \texttt{TradingState}, order submission.
\end{itemize}
\end{backgroundbox}

Each section uses a five-field template: Observable, Theory
chapter, Strategy skeleton, Pitfalls, Reference writeup.

%======================================================================
\section{Archetype 1: fixed-price market making}
\label{sec:archetype_mm}
%======================================================================

\subsection{Observable}

The product has an \emph{obvious, fixed true price} that does not
move.  Bids and asks cluster symmetrically around that level with a
spread wide relative to tick size, so every off-par fill yields
positive edge.  Takers periodically cross the book but the true price
never trends.

\begin{keyfactbox}[Diagnostic]
If you normalise price by the true value and see a stationary series
with zero drift, you have a fixed-price market-making product.
\end{keyfactbox}

\noindent
\textbf{Prosperity examples.}
\begin{itemize}
  \item \textbf{RAINFOREST\_RESIN} (Prosperity~3, Round~1): true price
  permanently fixed at 10{,}000~\citep{frankfurthedgehogs2025}.
  \item \textbf{AMETHYSTS} (Prosperity~2, Round~1): identical setup.
\end{itemize}

\subsection{Theory chapter}

\Cref{ch:market_making}, specifically Avellaneda--Stoikov.  Quote
bids below and asks above fair value, capture the spread on every
fill, manage inventory risk (one-sided accumulation) by skewing
quotes or flattening at the true price.

A known constant true price removes adverse-selection risk (no
informed counterparty) and fair-value estimation error.  The only
trade-off is edge versus fill probability: wider quotes earn more per
fill but fill less often.

\subsection{Strategy skeleton}

A complete market maker for a fixed-price product:

\begin{lstlisting}[caption={Fixed-price market maker (Prosperity).},
  label={lst:static_mm}]
TRUE_PRICE = 10000
EDGE = 2          # minimum distance from true price
LIMIT = 50        # position limit

def run(state, product):
    pos = state.position.get(product, 0)
    orders = []
    book = state.order_depths[product]
    # 1. Take any mispriced orders
    for ask, vol in sorted(book.sell_orders.items()):
        if ask < TRUE_PRICE and pos - vol <= LIMIT:
            orders.append(Order(product, ask, -vol))
            pos -= vol  # vol is negative for sells
    for bid, vol in sorted(book.buy_orders.items(),
                           reverse=True):
        if bid > TRUE_PRICE and pos - vol >= -LIMIT:
            orders.append(Order(product, bid, -vol))
            pos -= vol
    # 2. Quote passively
    buy_qty = min(LIMIT - pos, 20)
    sell_qty = min(LIMIT + pos, 20)
    if buy_qty > 0:
        orders.append(Order(product,
                            TRUE_PRICE - EDGE, buy_qty))
    if sell_qty > 0:
        orders.append(Order(product,
                            TRUE_PRICE + EDGE, -sell_qty))
    return orders
\end{lstlisting}

\noindent
Key design choices:
\begin{enumerate}
  \item \textbf{Take first, then make.}  Sweep any ask below true
  price or bid above it: free money.
  \item \textbf{Clip position.}  Bound \texttt{buy\_qty} and
  \texttt{sell\_qty} so you never breach the limit; exceeding it
  cancels \emph{all} orders (\cref{ch:prosperity_setup}).
  \item \textbf{Flatten at par.}  If inventory is maxed, suppress the
  passive quote on the overweight side.  Frankfurt Hedgehogs also
  flattened at exactly 10{,}000 to reset
  inventory~\citep{frankfurthedgehogs2025}.
\end{enumerate}

\subsection{Pitfalls}

\begin{pitfallbox}[Quoting too wide]
Too large an edge parameter means you never get filled, because other
market makers undercut you.  For RAINFOREST\_RESIN, Frankfurt Hedgehogs
overbid the best existing bid by one tick while keeping positive
edge; fill probability rose far more than edge decreased.
\end{pitfallbox}

\begin{pitfallbox}[Ignoring the position limit]
Prosperity does not partially reject orders.  If the sum of buy
quantities plus current position would exceed $+L$, \emph{every}
order submitted for that product is cancelled.  Cap order sizes
before submission.
\end{pitfallbox}

\subsection{Reference writeup}

Frankfurt Hedgehogs' \texttt{StaticTrader} used this pattern on
RAINFOREST\_RESIN, contributing $\approx$39{,}000 SeaShells per
round ($\approx$10\% of total Prosperity~3 PnL) from a strategy
``anyone could have come up with by carefully studying the
competition's matching rules''~\citep{frankfurthedgehogs2025}.

\begin{intuitionbox}
Fixed-price MM is the ``hello world'' of Prosperity: it teaches
take-then-make, position-limit clipping, and the edge-vs-fill
trade-off every later archetype inherits.
\end{intuitionbox}

\begin{prosperitybox}[title={RAINFOREST\_RESIN worked example}]
\textbf{Setup.}  True price $= 10{,}000$.  Position limit $= 50$.
Current position $= +12$.  The order book shows:

\medskip
\begin{tabular}{lrl}
\toprule
\textbf{Side} & \textbf{Price} & \textbf{Qty} \\
\midrule
Ask & 10{,}003 & 8 \\
Ask & 10{,}001 & 5 \\
Bid & 9{,}998  & 10 \\
Bid & 9{,}996  & 15 \\
\bottomrule
\end{tabular}
\medskip

\textbf{Step 1: Take.}  Ask at 10{,}001 is above true price, no bid
above 10{,}000, so there is nothing to sweep.

\textbf{Step 2: Make.}  Buy capacity $= 50 - 12 = 38$, sell capacity
$= 62$, both capped at clip size 20.  Post:
\begin{itemize}
  \item \textbf{Buy} 20 at 9{,}998 (matches bid wall; in practice
  overbid at 9{,}999 for queue priority).
  \item \textbf{Sell} 20 at 10{,}002 (undercuts ask at 10{,}003).
\end{itemize}
Both sides filling yields $2\times20 + 2\times20 = 80$ SeaShells per
timestep.  Across $\approx$10{,}000 timesteps, even modest fill rates
compound to tens of thousands.
\end{prosperitybox}

%======================================================================
\section{Archetype 2: trending / EMA / LR products}
\label{sec:archetype_trend}
%======================================================================

\subsection{Observable}

The price moves over time but follows a slow random walk with a
tight spread relative to tick size.  Normalising by Wall Mid
(\cref{ch:prosperity_setup}) makes the series look almost fixed-price.
Some products also show occasional sharp jumps or drift, creating a
second signal (trend-following or informed-flow detection).

\begin{keyfactbox}[Diagnostic]
Plot the raw series and its Wall Mid--normalised version.  If the
normalised series is stationary, use a moving-fair-value market
maker.  If the raw series shows persistent trends or jumps, overlay a
short EMA and check whether deviations predict reversals or
continuations.
\end{keyfactbox}

\noindent
\textbf{Prosperity examples.}
\begin{itemize}
  \item \textbf{KELP} (Prosperity~3, Round~1): slow random walk,
  tight spread.  Best treated as a dynamic-fair-value market maker.
  \item \textbf{SQUID\_INK} (Prosperity~3, Round~1): wider relative
  spread, sharp jumps.  Officially described as short-term
  mean-reverting, but the main edge came from an informed-trader signal
  (Olivia; see \cref{sec:archetype_signal}).
\end{itemize}

\subsection{Theory chapter}

Two chapters are relevant:
\begin{itemize}
  \item \Cref{ch:mean_reversion_ou}: if the normalised spread is
  stationary with mean-reverting tendencies, trade deviations from the
  EMA as described in the Ornstein--Uhlenbeck framework.
  \item \Cref{ch:momentum}: if the raw price shows persistent trends,
  a simple EMA crossover or linear-regression slope can generate
  trend-following signals.
\end{itemize}

The diagnostic question: \emph{mean-revert or trend?}  KELP was
``neither'': a random walk best traded as fixed-price MM with a
rolling fair-value estimate~\citep{frankfurthedgehogs2025}.
SQUID\_INK's description said ``mean-reverting,'' but no pure
price-based strategy was as reliable as detecting the informed trader
Olivia~\citep{frankfurthedgehogs2025}.

\subsection{Strategy skeleton}

\noindent
\textbf{Dynamic MM} (KELP-style).  Replace \texttt{TRUE\_PRICE} in
\cref{lst:static_mm} with a rolling \texttt{wall\_mid} (average of
bid-wall and ask-wall levels):

\begin{quote}
\texttt{fair = wall\_mid(book)}
\end{quote}

\noindent
Wall Mid is the midpoint between the deepest bid and ask levels
(ignoring small overbids/undercuts), giving a noise-robust estimate
of fair value.

\noindent
\textbf{EMA mean-reversion overlay} (SQUID\_INK pre-Olivia):

\begin{quote}
\texttt{ema = alpha * mid + (1 - alpha) * ema\_prev}\\
\texttt{if mid < ema - threshold: buy}\\
\texttt{if mid > ema + threshold: sell}
\end{quote}

\noindent
Frankfurt Hedgehogs used a fast EMA ($\alpha \approx 0.1$--$0.3$)
with fixed thresholds, avoiding z-score normalisation because they
saw no time-varying volatility~\citep{frankfurthedgehogs2025}.

\subsection{Pitfalls}

\begin{pitfallbox}[Confusing trend and mean-reversion]
A random walk looks like it ``trends'' short-horizon and
``mean-reverts'' long-horizon.  Without an autocorrelation test
(\cref{ch:mean_reversion_ou}), you overfit to whichever narrative you
looked for first.  KELP was near-pure random walk: both trend and
fade lost money after spread costs; MM around fair value won.
\end{pitfallbox}

\begin{pitfallbox}[Olivia bias]
Several Prosperity~3 products contained Olivia, an informed trader
who bought at daily lows and sold at daily highs.  Trading with
Olivia on SQUID\_INK was worth $\approx$8{,}000 SeaShells per round;
ignoring her and running generic mean-reversion yielded high
variance, near-zero expected PnL.  Check for informed-flow before
defaulting to price-only models
(\cref{sec:archetype_signal}).
\end{pitfallbox}

\subsection{Reference writeup}

Frankfurt Hedgehogs' \texttt{DynamicTrader} traded KELP as moving
fair-value MM; their \texttt{InkTrader} traded SQUID\_INK purely on
Olivia's daily-extrema pattern, with no MM or mean-reversion for that
product~\citep{frankfurthedgehogs2025}.

\begin{intuitionbox}
Archetype~2's lesson: disciplined diagnosis.  Before throwing an EMA
at a moving price, ask: ``Is there a predictable signal, or is this a
random walk?''  If random walk, MM around fair value; if signal,
identify its source (price vs.\ informed-flow) before choosing.
\end{intuitionbox}

\begin{prosperitybox}[title={KELP: Wall Mid vs.\ raw mid}]
At a given timestep, KELP's order book shows:

\medskip
\begin{tabular}{lrr}
\toprule
\textbf{Side} & \textbf{Price} & \textbf{Qty} \\
\midrule
Ask & 2{,}032 & 3 \\
Ask & 2{,}031 & 22 \quad$\leftarrow$ ask wall \\
Bid & 2{,}029 & 18 \quad$\leftarrow$ bid wall \\
Bid & 2{,}027 & 5 \\
\bottomrule
\end{tabular}
\medskip

Raw mid $= (2{,}031 + 2{,}029)/2 = 2{,}030$.  The ask at 2{,}032
comes from a small overbidding bot and distorts the raw mid.  Wall
Mid (deepest-liquidity levels) is also 2{,}030 here, but if a bot
undercuts the ask wall with 3 lots at 2{,}028, the raw mid drops to
2{,}028.5 while Wall Mid stays at 2{,}030.

Wall Mid filters noise from small overbids/undercuts, stabilising the
fair-value estimate.  Frankfurt Hedgehogs found this ``crucial for
designing effective strategies'' across all
products~\citep{frankfurthedgehogs2025}.
\end{prosperitybox}

%======================================================================
\section{Archetype 3: basket / ETF statistical arbitrage}
\label{sec:archetype_basket}
%======================================================================

\subsection{Observable}

One or two \emph{baskets} (ETF analogues) plus several constituents
appear together.  Each basket is a known weighted combination of
constituents, trades on its own book, and has \emph{no
creation/redemption} mechanism.  The spread between basket and
synthetic (weighted constituent sum) mean-reverts around a non-zero
premium.

\begin{keyfactbox}[Diagnostic]
Compute $\text{spread}_t = P_{\text{basket},t} - \sum_i w_i \,
P_{i,t}$.  If this spread is stationary and mean-reverting, you have
a basket stat-arb product.
\end{keyfactbox}

\noindent
\textbf{Prosperity examples.}
\begin{itemize}
  \item \textbf{PICNIC\_BASKET1} (Prosperity~3, Round~2): composition
  $= 6 \times \text{CROISSANTS} + 3 \times \text{JAMS} + 1 \times
  \text{DJEMBES}$.
  \item \textbf{PICNIC\_BASKET2} (Prosperity~3, Round~2): composition
  $= 4 \times \text{CROISSANTS} + 2 \times \text{JAMS}$.
  \item Constituent products: CROISSANTS, JAMS, DJEMBES.
\end{itemize}

\subsection{Theory chapter}

\Cref{ch:pairs_cointegration} (cointegration) and
\cref{ch:basket_etf_arb} (ETF-specific).  The basket--synthetic
spread is analogous to real-world ETF
premium/discount~\citep{hull8ed}; Prosperity's absence of
creation/redemption lets the spread stay wider and more persistent.

\subsection{Strategy skeleton}

Core loop:

\begin{quote}
\texttt{synthetic = 6*mid(CROISSANTS) + 3*mid(JAMS) + 1*mid(DJEMBES)}\\
\texttt{spread = mid(BASKET1) - synthetic - premium}\\
\texttt{if spread > +threshold: sell basket, (optionally) buy legs}\\
\texttt{if spread < -threshold: buy basket, (optionally) sell legs}\\
\texttt{if spread crosses zero: flatten}
\end{quote}

\noindent
\textbf{Key design choices:}
\begin{enumerate}
  \item \textbf{Subtract the running premium.}  The spread mean is
  non-zero: baskets trade at a persistent premium.  Frankfurt
  Hedgehogs tracked and subtracted a running
  estimate~\citep{frankfurthedgehogs2025}.
  \item \textbf{Fixed thresholds.}  A light grid search prioritising
  ``landscape stability'' (flat regions of good PnL) over peak
  performance~\citep{frankfurthedgehogs2025}.
  \item \textbf{Half-hedge.}  In the final round they hedged 50\% of
  basket exposure, a compromise between full-hedge variance
  reduction and unhedged expected value.
  \item \textbf{Flatten at zero.}  Close on zero-crossing (adjusted
  for Olivia) rather than waiting for the opposite threshold.
\end{enumerate}

\subsection{Pitfalls}

\begin{pitfallbox}[Non-zero premium]
Assuming zero-mean reversion when the spread actually reverts to a
persistent premium biases every entry.  Compute the running spread
average and subtract it.  Frankfurt Hedgehogs updated this premium
continuously in live trading.
\end{pitfallbox}

\begin{pitfallbox}[Olivia cross-product signal]
Olivia appeared on CROISSANTS (in both baskets), not the baskets
themselves.  Olivia-short Croissants meant the basket spread would
widen.  Frankfurt Hedgehogs biased basket thresholds accordingly: if
Olivia short, shift long entry $-50 \to -80$, short entry $+50 \to
+20$~\citep{frankfurthedgehogs2025}.  Teams ignoring this
cross-product signal left 20{,}000+ SeaShells per round on the table.
\end{pitfallbox}

\begin{pitfallbox}[No creation/redemption]
Real-world ETF arb has authorised participants creating/redeeming
shares to enforce NAV tracking.  Prosperity has no such mechanism;
basket and constituents are merely correlated, so spreads persist
longer and move further than textbook ETF arb predicts.
\end{pitfallbox}

\subsection{Reference writeup}

Frankfurt Hedgehogs' \texttt{EtfTrader} implemented the
spread-threshold model with Olivia-adjusted entries; combined basket
and CROISSANTS trading generated $\approx$60{,}000--80{,}000
SeaShells per round.  The chrispyroberts writeup complements this
with a systematic grid search over threshold
parameters\footnote{See \texttt{research/writeups/chrispyroberts\_P3/}.}.

\begin{intuitionbox}
Basket stat-arb in Prosperity is the closest analogue to real-world
ETF arbitrage absent creation/redemption.  That absence is both the
opportunity (wider spreads, more PnL) and the risk (longer
dislocations).  Winning recipe: fixed thresholds, premium adjustment,
and, if present, an informed-flow signal on a constituent.
\end{intuitionbox}

\begin{prosperitybox}[title={PICNIC\_BASKET1 spread calculation}]
Mid prices:

\medskip
\begin{tabular}{lr}
\toprule
\textbf{Product} & \textbf{Mid price} \\
\midrule
CROISSANTS & 4{,}330 \\
JAMS       & 6{,}540 \\
DJEMBES    & 9{,}800 \\
PICNIC\_BASKET1 & 56{,}200 \\
\bottomrule
\end{tabular}
\medskip

Synthetic $= 6\cdot 4{,}330 + 3\cdot 6{,}540 + 1\cdot 9{,}800 =
55{,}400$.  Raw spread $= 56{,}200 - 55{,}400 = +800$.  With running
premium $+400$, adjusted spread $= +400$.

Entry threshold $\pm 300$: we are above $+300$, so the basket is
``expensive.''  \textbf{Action:} sell PICNIC\_BASKET1, optionally
half-hedge ($3\times$CROISSANTS, $1.5\times$JAMS, $0.5\times$DJEMBES,
rounded).  Hold until adjusted spread crosses zero; flatten.

If Olivia was short CROISSANTS (falling prices expected, spread
widening), thresholds shifted to $\{-530, +70\}$: harder longs,
easier shorts~\citep{frankfurthedgehogs2025}.
\end{prosperitybox}

%======================================================================
\section{Archetype 4: options / Black--Scholes + IV hedging}
\label{sec:archetype_options}
%======================================================================

\subsection{Observable}

An underlying plus several \emph{vouchers} (call options) at
different strikes, with time-to-expiry shrinking each round.  Implied
volatility (IV) vs.\ moneyness forms a parabolic smile; option
prices fluctuate rapidly around Black--Scholes theoretical values,
creating short-term mean-reversion opportunities.

\begin{keyfactbox}[Diagnostic]
Compute each strike's IV via Black--Scholes inversion
(\cref{ch:black_scholes}) and plot against moneyness.  If points lie
on a smooth parabola, fit it, detrend, and trade the
mean-reverting residuals.
\end{keyfactbox}

\noindent
\textbf{Prosperity examples.}
\begin{itemize}
  \item \textbf{VOLCANIC\_ROCK\_VOUCHERS} (Prosperity~3, Round~3):
  five call options on VOLCANIC\_ROCK at strikes 9{,}500, 9{,}750,
  10{,}000, 10{,}250, and 10{,}500.  Expiry shrank from 7 days
  (Round~3) to 2 days (Round~5).
  \item Underlying: \textbf{VOLCANIC\_ROCK}, trading around 10{,}000.
\end{itemize}

\subsection{Theory chapter}

Three chapters:
\begin{itemize}
  \item \Cref{ch:black_scholes}: Black--Scholes, risk-neutral
  valuation, IV inversion.
  \item \Cref{ch:greeks_hedging}: delta, gamma, vega, theta,
  portfolio delta-hedging.
  \item \Cref{ch:vol_surface}: vol smile, parametric surface fits,
  smile-dynamics relative value.
\end{itemize}

Frankfurt Hedgehogs combined two alpha
sources~\citep{frankfurthedgehogs2025}:

\begin{enumerate}
  \item \textbf{IV scalping.}  Fit a parabola to IV vs.\ moneyness,
  convert fitted IV back to a fair price via Black--Scholes, trade
  the residual.  Analogous to real-desk relative-value
  vol~\citep{sinclair2013volatility}.
  \item \textbf{Gamma scalping.}  Long gamma plus delta rehedging:
  convexity profit exceeds theta when realised volatility beats
  implied.
\end{enumerate}

\subsection{Strategy skeleton}

\begin{quote}
\texttt{for each strike K:}\\
\quad\texttt{iv[K] = bs\_invert(market\_price[K], S, K, T, r)}\\
\texttt{fit parabola: iv\_hat = a*m\^{}2 + b*m + c}\\
\quad\texttt{where m = log(S/K) / sqrt(T)}\\
\quad\texttt{\# m is moneyness: how far the strike K sits above/below spot S}\\
\texttt{for each strike K:}\\
\quad\texttt{fair\_price = bs\_call(S, K, T, r, iv\_hat[K])}\\
\quad\texttt{deviation = market\_price[K] - fair\_price}\\
\quad\texttt{if deviation > +threshold: sell voucher K}\\
\quad\texttt{if deviation < -threshold: buy voucher K}
\end{quote}

\noindent
Frankfurt Hedgehogs started with the 10{,}000 (ATM) strike for
liquidity and expanded to others as the underlying moved or expiry
shortened~\citep{frankfurthedgehogs2025}.

\subsection{Pitfalls}

\begin{pitfallbox}[Wrong volatility estimate]
A poorly calibrated parabola (too few points, or outliers from deep
OTM options with near-zero extrinsic value) yields a bad ``fair
price'' and trades noise.  Discard outliers where extrinsic value is
too low for reliable IV.
\end{pitfallbox}

\begin{pitfallbox}[Discrete delta-hedging costs]
Continuous delta-hedging is free in theory; in Prosperity (integer
quantities, discrete timesteps, bid--ask spread on the underlying),
each rebalance costs spread.  Frankfurt Hedgehogs found explicit
delta-hedging ``prohibitively expensive'' and used a hybrid: IV
scalping primary, a moderate mean-reversion position in the
underlying as luck
hedge~\citep{frankfurthedgehogs2025}.
\end{pitfallbox}

\begin{pitfallbox}[Integer position constraints]
Option positions are integers: delta 0.53 cannot be hedged exactly.
Rounding errors accumulate across multiple strikes.  Solution:
aggregate portfolio delta, hedge at portfolio level.
\end{pitfallbox}

\subsection{Reference writeup}

Frankfurt Hedgehogs' \texttt{OptionTrader} did IV scalping with a
parabolic smile fit; it contributed $\approx$100{,}000--150{,}000
SeaShells per round, their single most profitable
strategy~\citep{frankfurthedgehogs2025}.  A secondary VOLCANIC\_ROCK
mean-reversion position produced volatile returns ($+$100k, $-$50k,
$-$10k across rounds), serving as a relative hedge against other top
teams pursuing pure mean-reversion.

\begin{intuitionbox}
Prosperity options are \emph{relative value on the vol smile}, not
directional bets.  Fit the smile, find the
outlier, trade it back.  Exactly what a real vol desk does
(\cref{ch:vol_surface}); Prosperity just enlarges the mispricings and
speeds up feedback.
\end{intuitionbox}

\begin{prosperitybox}[title={IV scalping walk-through}]
\textbf{Setup.}  VOLCANIC\_ROCK trades at $S = 10{,}050$.  Time to
expiry $T = 5$ days $= 5/365 \approx 0.0137$ years.  Risk-free rate
$r = 0$.  We observe market prices for five voucher strikes:

\medskip
\begin{tabular}{rrrrr}
\toprule
\textbf{Strike} & \textbf{Market price} & \textbf{Moneyness $m$} &
\textbf{IV (BS invert)} & \textbf{IV (fitted)} \\
\midrule
9{,}500  & 565 & $+0.449$ & 0.162 & 0.160 \\
9{,}750  & 340 & $+0.247$ & 0.155 & 0.152 \\
10{,}000 & 145 & $+0.041$ & 0.148 & 0.150 \\
10{,}250 &  38 & $-0.164$ & 0.160 & 0.155 \\
10{,}500 &   6 & $-0.368$ & 0.175 & 0.168 \\
\bottomrule
\end{tabular}
\medskip

Moneyness $m = \ln(S/K) / \sqrt{T}$.  Fit parabola $\hat{v}(m) = am^2
+ bm + c$ to the five IVs via least squares; the fitted column is
the smooth smile.

\textbf{Signal.}  At the 10{,}250 strike, market IV $= 0.160$,
fitted $= 0.155$.  $\text{BS}(S, 10250, T, r, 0.155) = 34$; market
price 38 is overpriced by 4 SeaShells.

\textbf{Action.}  Sell at 38, buy back near fitted 34.  Expected
edge $\approx 4$ SeaShells per lot; repeated across strikes and
timesteps, this compounds to the 100k+ Frankfurt Hedgehogs reported.
\end{prosperitybox}

%======================================================================
\section{Archetype 5: cross-exchange / locational arbitrage}
\label{sec:archetype_location}
%======================================================================

\subsection{Observable}

The product trades locally \emph{and} externally (separate venue
with fixed bid/ask quotes), but cross-venue transfer incurs
import/export fees, transport, tariffs, or storage charges.  The
external price is observable only through a ``conversion'' mechanism
with per-timestep capacity.

\begin{keyfactbox}[Diagnostic]
If a product has both an order book and a conversion mechanism,
compute the round-trip all-in cost.  If it is sometimes negative,
locational arbitrage exists.
\end{keyfactbox}

\noindent
\textbf{Prosperity examples.}
\begin{itemize}
  \item \textbf{MAGNIFICENT\_MACARONS} (Prosperity~3, Round~4):
  local-exchange tradable plus external conversion.  Conversion limit
  10/timestep, position limit 75.  External price depended on
  sunlight, sugar, shipping, tariffs~\citep{frankfurthedgehogs2025}.
\end{itemize}

\subsection{Theory chapter}

\Cref{ch:arbitrage_zoo}, locational-arbitrage section.  Classic
rule: if $\bidp_{\text{local}} > \askp_{\text{external}} + \text{fees}$,
sell local and buy external (or vice versa).  Prosperity twists:
limited conversion capacity and a hidden taker bot that fills above
the visible book.

\subsection{Strategy skeleton}

\begin{quote}
\texttt{ext\_ask\_cost = external\_ask + import\_tariff + transport}\\
\texttt{ext\_bid\_revenue = external\_bid - export\_tariff - transport}\\
\texttt{if best\_bid\_local > ext\_ask\_cost:}\\
\quad\texttt{sell local, buy external (convert +10)}\\
\texttt{if best\_ask\_local < ext\_bid\_revenue:}\\
\quad\texttt{buy local, sell external (convert -10)}\\
\texttt{\# Hidden taker exploitation:}\\
\texttt{offer\_price = int(ext\_bid\_revenue + 0.5)}\\
\texttt{place limit sell at offer\_price for 10 units}\\
\texttt{if filled: convert -10 to close position}
\end{quote}

\noindent
The last three lines capture Frankfurt Hedgehogs' insight: a hidden
taker fills sells at \texttt{int(externalBid + 0.5)} with
$\approx$60\% probability, $\approx$3 SeaShells edge per unit over
naive best-bid~\citep{frankfurthedgehogs2025}.

\subsection{Pitfalls}

\begin{pitfallbox}[Asymmetric conversion costs]
Import/export fees are typically asymmetric (e.g.\ import 10, export
5), so arbitrage thresholds differ by direction.  Compute all-in cost
separately per leg.
\end{pitfallbox}

\begin{pitfallbox}[Transport / conversion delay]
Conversions reduce position but settle with a one-timestep lag.
Converting 10 units consumes that capacity for the next timestep.
Size quotes to conversion limit, not position limit.
\end{pitfallbox}

\subsection{Reference writeup}

Frankfurt Hedgehogs' \texttt{CommodityTrader} handled
MAGNIFICENT\_MACARONS: limit sells at the hidden-taker price,
convert surplus externally.  Despite conservative quoting (10 units
per timestep), the strategy generated $\approx$80{,}000--100{,}000
SeaShells per round.  They noted quoting 20--30 units would have
been more profitable since excess inventory still converts
profitably~\citep{frankfurthedgehogs2025}.

A similar hidden-taker structure appeared with \textbf{ORCHIDS} in
Prosperity~2 (LinearUtility
writeup\footnote{See \texttt{research/writeups/LinearUtility\_P2/}.}).
Studying prior-year writeups was a direct competitive advantage.

\begin{intuitionbox}
Prosperity locational arb goes beyond the pure textbook version: it
requires discovering hidden structure (the taker bot) empirically.
\Cref{ch:arbitrage_zoo} tells you \emph{where} to look;
competition-specific microstructure tells you \emph{how much}.
\end{intuitionbox}

\begin{prosperitybox}[title={MAGNIFICENT\_MACARONS profit arithmetic}]
External bid (net of fees) nets 640; local best bid is 638.  Naive
arb says ``no opportunity.''  But offering to sell locally at
$\lfloor 640 + 0.5 \rfloor = 640$ and being filled by the hidden
taker $\approx$60\% of the time: sell at 640, convert at all-in cost
637, for 3 SeaShells edge per filled unit.

With 10 units per timestep $\times$ 10{,}000 timesteps:
\[
  \text{Theoretical max PnL} = 300{,}000 \;\text{SeaShells.}
\]
At 60\% fill, expected PnL $\approx 180{,}000$.  Frankfurt Hedgehogs
realised $\approx$80{,}000--100{,}000 because edge was not always 3
and they quoted conservatively; quoting 20--30 units would have
pushed PnL to $\approx$130{,}000--160{,}000.
\end{prosperitybox}

%======================================================================
\section{Archetype 6: signal-driven / informed-flow exploitation}
\label{sec:archetype_signal}
%======================================================================

\subsection{Observable}

One or more bot traders exhibit predictable patterns tied to price
extremes, time-of-day, or other detectable features.  Prosperity~3's
key case was \textbf{Olivia}: bought 15 lots at daily lows, sold 15
at daily highs across SQUID\_INK and CROISSANTS.  Detectable even
before trader IDs appeared in Round~5.

\begin{keyfactbox}[Diagnostic]
Filter trades by quantity (e.g.\ exactly 15 lots).  If trades of
that size cluster at daily extremes and align with subsequent
reversals, you have an informed-flow signal.
\end{keyfactbox}

\noindent
\textbf{Prosperity examples.}
\begin{itemize}
  \item \textbf{SQUID\_INK} (Prosperity~3, Rounds 1--5): Olivia bought
  15 lots at daily lows, sold 15 lots at daily highs.
  \item \textbf{CROISSANTS} (Prosperity~3, Rounds 2--5): same Olivia
  pattern, exploitable directly \emph{and} as a cross-product signal
  for the baskets.
\end{itemize}

\subsection{Theory chapter}

Three chapters give the backdrop:
\begin{itemize}
  \item \Cref{ch:kyle_gm}: Glosten--Milgrom.  Olivia is an
  \emph{informed insider} whose trades reveal private information on
  price direction.
  \item \Cref{ch:order_flow}: order-flow analysis, toxic-flow
  detection, VPIN-style metrics.
  \item \Cref{ch:ml_trading}: feature engineering on trade data
  (quantity, timing, direction relative to extremes).
\end{itemize}

\subsection{Strategy skeleton}

\begin{quote}
\texttt{daily\_low = running\_min(mid\_prices\_today)}\\
\texttt{daily\_high = running\_max(mid\_prices\_today)}\\
\texttt{for each trade in recent\_trades:}\\
\quad\texttt{if trade.qty == 15 and trade.price <= daily\_low:}\\
\qquad\texttt{signal = LONG  \# Olivia bought at low; price will rise}\\
\quad\texttt{if trade.qty == 15 and trade.price >= daily\_high:}\\
\qquad\texttt{signal = SHORT \# Olivia sold at high; price will fall}\\
\texttt{position toward signal; flatten on new extreme or EOD}
\end{quote}

\noindent
Once trader IDs appeared in Round~5, detection simplified to
\texttt{if trade.buyer == ``Olivia''}, removing false positives from
other 15-lot trades~\citep{frankfurthedgehogs2025}.

\subsection{Pitfalls}

\begin{pitfallbox}[Direction flips]
A new daily extreme invalidates Olivia's signal: if she bought at
the then-low but the price makes a new low, the original signal was
a false positive.  Monitor for new extrema and reset.
\end{pitfallbox}

\begin{pitfallbox}[Detection latency before trader IDs]
Without trader IDs (Rounds 1--4), detection relied on 15-lot
quantity plus price-extreme matching.  Other bots occasionally hit
15 lots at non-extreme prices, producing false positives with low
per-event cost but lower Sharpe.
\end{pitfallbox}

\subsection{Reference writeup}

Frankfurt Hedgehogs detected Olivia early and exploited the signal
on SQUID\_INK ($\approx$8{,}000/round) and CROISSANTS
($\approx$20{,}000/round).  The CROISSANTS signal also fed the
basket thresholds in Archetype~3 (\cref{sec:archetype_basket}),
amplifying PnL~\citep{frankfurthedgehogs2025}.  Identifiable-bot
exploits are not new: similar patterns appeared in Prosperity~1--2;
teams studying prior writeups (Stanford Cardinal, Jasper) entered
Prosperity~3 with detection code ready.

\begin{prosperitybox}[title={Detecting Olivia on SQUID\_INK}]
Today's SQUID\_INK mid range: low 1{,}820 (timestep 340), high
1{,}895 (timestep 680).  At timestep 700 the tape shows a new trade:

\medskip
\begin{tabular}{lrrll}
\toprule
\textbf{Timestep} & \textbf{Price} & \textbf{Qty} & \textbf{Buyer} &
\textbf{Seller} \\
\midrule
700 & 1{,}895 & 15 & --- & \textit{(unknown)} \\
\bottomrule
\end{tabular}
\medskip

A 15-lot sell at the daily high.  Without trader IDs, flag as a
candidate Olivia sell on two criteria: quantity $= 15$ and price
$\geq$ daily high.  Prediction: SQUID\_INK falls from here.

\textbf{Action:} go short to position limit.  Hold until (a)~a 15-lot
buy at a new daily low (Olivia reversal), (b)~a new daily high above
1{,}895 (invalidation), or (c)~end of day.

\textbf{Invalidation:} if SQUID\_INK hits 1{,}900 at timestep 750,
the original sell was not the true maximum; flatten immediately.

This detection-and-invalidation logic kept the strategy robust:
false positives closed quickly before losses compounded.
\end{prosperitybox}

\begin{intuitionbox}
Olivia is a simplified informed-trading model
(\cref{ch:kyle_gm}).  Real markets hide individual traders, but
the \emph{pattern} (unusual sizes at unusual prices preceding
reversals) is exactly what VPIN and order-flow-toxicity metrics
target.  Prosperity makes it concrete.
\end{intuitionbox}

%======================================================================
\section{Archetype-recognition flowchart}
\label{sec:archetype_flowchart}
%======================================================================

Run new products through the following decision table, which is the
Prosperity-specific version of \cref{ch:which_strategy_when}'s
framework.

\begin{table}[htbp]
\centering
\small
\caption{Archetype recognition: observable features $\to$ archetype
  $\to$ strategy.}
\label{tab:archetype_recognition}
\begin{tabular}{@{}p{4.2cm}lp{4.5cm}@{}}
\toprule
\textbf{Observable feature} & \textbf{Archetype} &
\textbf{Primary strategy} \\
\midrule
Fixed known true price; wide spread &
1 (\cref{sec:archetype_mm}) &
Fixed-price market making \\
\addlinespace
Moving price; tight spread; no trend &
2 (\cref{sec:archetype_trend}) &
Dynamic market making (Wall Mid) \\
\addlinespace
Moving price; sharp jumps; 15-lot trades at extremes &
6 (\cref{sec:archetype_signal}) &
Informed-flow detection (Olivia) \\
\addlinespace
Multiple products; known basket composition; mean-reverting spread &
3 (\cref{sec:archetype_basket}) &
Basket stat-arb, threshold entry \\
\addlinespace
Options (vouchers) on an underlying; parabolic IV smile &
4 (\cref{sec:archetype_options}) &
IV scalping + optional gamma scalp \\
\addlinespace
Local + external exchange; conversion fees &
5 (\cref{sec:archetype_location}) &
Locational arb + hidden-taker exploit \\
\bottomrule
\end{tabular}
\end{table}

\noindent
The table orders archetypes by diagnostic speed.  Fixed-price
products are obvious in seconds; informed-flow signals take hours of
trade-data study.  Time budget per round:

\begin{enumerate}
  \item \textbf{Minutes 0--5.}  Test for fixed price (Archetype~1) or
  conversion mechanism (Archetype~5); both are instant to spot and
  fast to code.
  \item \textbf{Minutes 5--30.}  Check for baskets/ETFs
  (Archetype~3); compute spread, test stationarity.
  \item \textbf{Minutes 30--120.}  For options (Archetype~4): IVs,
  smile fit, mean-reverting residuals.
  \item \textbf{Hours 2--6.}  Remaining products: trade-data analysis
  for informed flow (Archetype~6), autocorrelation tests
  (Archetype~2), normalised order-book plots.
\end{enumerate}

\begin{keyfactbox}[The five-minute rule]
A product unclassifiable within five minutes is Archetype~2
(trending/EMA) or Archetype~6 (signal-driven).  Default to dynamic
MM while investigating; it extracts spread PnL without a directional
commitment.
\end{keyfactbox}

\subsection{PnL breakdown by archetype}

\Cref{tab:pnl_archetype} lists Frankfurt Hedgehogs' approximate
per-round PnL by archetype.  The lesson: the simplest strategies
(fixed-price MM, locational arb) delivered the most reliable
returns; the most sophisticated (IV scalping) had the highest
absolute PnL but more variance.

\begin{table}[htbp]
\centering
\small
\caption{Approximate PnL per round by archetype (Frankfurt Hedgehogs,
  Prosperity~3).  Numbers are in SeaShells.}
\label{tab:pnl_archetype}
\begin{tabular}{@{}llr@{}}
\toprule
\textbf{Archetype} & \textbf{Products} &
\textbf{PnL / round} \\
\midrule
1: Fixed-price MM & RAINFOREST\_RESIN & $\approx$39{,}000 \\
2: Dynamic MM & KELP & $\approx$5{,}000 \\
2/6: Trend + signal & SQUID\_INK & $\approx$8{,}000 \\
3: Basket stat-arb & PICNIC\_BASKET1/2 + CROISSANTS &
  $\approx$60{,}000--80{,}000 \\
4: Options IV scalp & VOLCANIC\_ROCK\_VOUCHERS &
  $\approx$100{,}000--150{,}000 \\
5: Locational arb & MAGNIFICENT\_MACARONS &
  $\approx$80{,}000--100{,}000 \\
\midrule
\textbf{Total} & & $\approx$\textbf{290{,}000--380{,}000} \\
\bottomrule
\end{tabular}
\end{table}

\begin{historybox}
\textbf{Archetypes recur across Prosperity editions.}  P1 (2023)
introduced fixed-price MM (Amethysts), trend (Starfruit), basket arb
(Gift Baskets), locational (Orchids).  P2 (2024) repeated the same
archetypes with renamed products.  P3 (2025) added options (Volcanic
Rock Vouchers) and foregrounded Olivia.  Teams who studied
prior-year writeups (Stanford Cardinal in P1, Jasper/LinearUtility
in P2, Frankfurt Hedgehogs in P3) entered each competition with
detection code, parameter ranges, and strategy skeletons ready.
Highest-Sharpe pre-competition activity available.
\end{historybox}

%======================================================================
\section*{}
\vspace{-2em}
\begin{gsbox}
\textbf{New Strat on day one.}  A new Goldman Strat given an
unfamiliar product does the same thing you do in Prosperity:
classify.  Flow product (desk market-makes)?  Relative value (trade
basis vs.\ hedge)?  Vol product (scalp the smile)?  Cross-venue
(arb exchange vs.\ OTC)?  The classification picks the model
library, risk limits, and P\&L buckets.  Prosperity compresses this
diagnostic into a 48-hour round.
\end{gsbox}

%======================================================================
\section{Key facts to memorise}
\label{sec:algo_rounds_keyfacts}
%======================================================================

\begin{enumerate}
  \item \textbf{Six archetypes} cover every Prosperity product to
  date: fixed-price MM, trending/EMA, basket stat-arb, options/IV
  scalping, locational arb, informed-flow exploitation.

  \item \textbf{Fixed-price} (RAINFOREST\_RESIN, AMETHYSTS): known
  constant price; quote around it.  Simplest archetype; implement
  within 30 minutes of identifying.

  \item \textbf{Dynamic} (KELP): random walks.  Replace fixed price
  with Wall Mid; otherwise identical to Archetype~1.

  \item \textbf{Baskets} (PICNIC\_BASKET1/2): mean-revert to
  synthetic value \emph{plus persistent premium}.  Fixed thresholds
  (not z-scores), subtract premium, half-hedge if unsure.

  \item \textbf{Options} (VOLCANIC\_ROCK\_VOUCHERS): relative value
  on the IV smile.  Fit parabola, trade residuals, hedge at
  portfolio level.  IV scalping was P3's single highest-PnL strategy
  ($\approx$100k--150k/round).

  \item \textbf{Locational arb} (MAGNIFICENT\_MACARONS): textbook
  arb plus hidden taker-bot discovery.  Similar structures recur;
  study prior writeups.

  \item \textbf{Informed-flow} (Olivia on SQUID\_INK, CROISSANTS):
  low-risk directional bets.  Detect by quantity + price-extreme;
  use cross-product signals to bias basket thresholds.

  \item \textbf{Take first, then make.}  Sweep mispriced orders
  before placing passive quotes.

  \item \textbf{Clip every order} against the position limit;
  exceeding it cancels \emph{all} orders.

  \item \textbf{Diagnose before coding.}  Spend the first 30 minutes
  classifying products.  \Cref{tab:archetype_recognition} plus the
  time-allocation schedule give a checklist.

  \item \textbf{Study prior-year writeups.}  Every Frankfurt
  Hedgehogs P3 alpha source had a P1/P2 predecessor.  Highest-Sharpe
  pre-competition activity.

  \item \textbf{Robustness over peak backtest PnL.}  Grid search for
  flat regions, keep parameter counts low.  Unexplainable backtest
  performance is probably noise.
\end{enumerate}
